\chapter{Stress-Energy and the Einstein Field Equations}

These notes follow Leonard Susskind's ninth lecture on general relativity, with transcript and board evidence curated with LazyingArt LLC and Video2Book. The lecture proceeds in a deliberate way. We are not going to win by brute force tensor algebra; we are going to understand what the field equation must look like, why the source cannot be only a mass density, and why a very particular combination of curvature terms is forced on us. We set \(c=1\) throughout.

\section{Why We Will Not Solve the Equations Head-On}

Susskind opens by lowering the temperature. Einstein's equations are not conceptually obscure in the way one first fears; what is obscure is how quickly the calculations become nasty. Christoffel symbols multiply, curvature components multiply again, and every object seems to carry derivatives of other objects. Even writing the equations down in components is already unpleasant.

That is why the lecture does not try to solve the field equations in full generality. If we really want to compute Christoffel symbols, Ricci tensors, or Einstein tensors for a metric, a computer is the natural ally. Mathematica can do the heavy lifting. The point of the lecture is different: we want the meaning of the symbols, the logic of the source term, and the criterion that selects the correct geometric tensor.

So the lecture backs up to Newtonian gravity. There the same two-way logic is already present in a much simpler form.

\section{Newton's Two-Way Street and a Heuristic Bridge}

Newtonian gravity already has the right architecture. The field tells matter how to move, and matter tells the field what to be. The motion law is
\begin{align}
m\mathbf a &= -m\nabla\phi,\\
\mathbf a &= -\nabla\phi,
\end{align}
where \(\phi\) is the gravitational potential.

The field equation is Poisson's equation,
\begin{equation}
\nabla^2\phi = 4\pi G\rho,
\end{equation}
with \(\rho\) the mass density. For a spherically symmetric source of total mass \(M\), the exterior solution is
\begin{equation}
\phi(r)=-\frac{GM}{r}.
\end{equation}
Outside the source this solves the field equation, and taking the gradient reproduces the familiar inverse-square force law.

This is the two-way street in its simplest form: \(\phi\) determines acceleration, while \(\rho\) determines \(\phi\).

Now we ask what this structure ought to become in general relativity. Susskind does not write down the full Schwarzschild metric again. He only recalls the time-time component,
\begin{equation}
g_{00}=1-\frac{2GM}{r}.
\end{equation}
Since \(\phi=-GM/r\), we can rewrite this as
\begin{equation}
g_{00}=1+2\phi.
\end{equation}
Therefore
\begin{equation}
\nabla^2 g_{00}=2\nabla^2\phi\sim 8\pi G\rho.
\end{equation}

This must be taken with a grain of salt. Susskind says explicitly that it is only a mnemonic. It is not Einstein's equation. But it is a good mnemonic, because it already has the rough texture we expect: matter appears to be sourcing geometry. The old Newtonian law \(\mathbf a=-\nabla\phi\) will be replaced by the statement that particles move on spacetime geodesics, and the crude source equation above will be replaced by the real field equation.

Before we can write that real equation down, however, we need to understand the right-hand side. Mass means energy, and energy density is not a lone scalar. It is part of a larger relativistic object.

\section{From Charge Conservation to the Energy-Momentum Tensor}

Mass really means energy. With \(c=1\), \(E=mc^2\) simply becomes the identification of mass and energy, so what looked Newtonianly like a mass density should already be read as an energy density. But energy density is not enough. It belongs to a whole complex of quantities, and that is why the lecture backs up and reviews local conservation.

Susskind starts with electric charge because charge is simpler than energy and momentum. Let \(Q\) be the total charge of a system, let \(\sigma\) be the charge density, and let \(\mathbf J\) be the spatial current of charge, the amount of charge crossing an oriented infinitesimal window per unit area per unit time. Local conservation means that charge cannot disappear here and reappear at Alpha Centauri. If charge leaves a small box, it must pass through the walls. Thus
\begin{equation}
\dot{\sigma}+\nabla\cdot\mathbf J=0.
\end{equation}

If we package density and current into the four-current
\begin{equation}
J^\mu=(J^0,J^1,J^2,J^3), \qquad J^0=\sigma,
\end{equation}
then the continuity equation becomes
\begin{equation}
\partial_\mu J^\mu=0.
\end{equation}
In flat spacetime this is the clean form. In curved coordinates one promotes ordinary derivatives to covariant derivatives so that the equation keeps its tensorial meaning.

This example teaches two things at once. First, density and flow belong together. Second, good laws of physics should be tensor laws, valid in any coordinates.

Now we promote the same logic from charge to energy and momentum. Here there is one extra twist. Electric charge itself is invariant, whereas total energy and total momentum are not. They are the components of the four-vector
\begin{equation}
P^\mu=(P^0,P^1,P^2,P^3), \qquad P^0=E.
\end{equation}
Each component of \(P^\mu\) is a conserved quantity of the same general type as charge: it has a density and it has a flow. That is why one more index is needed.

\paragraph{What the two indices mean.} The first index tells us which conserved quantity we are talking about. The value \(0\) means energy; the values \(1,2,3\) mean the three momentum components. The second index tells us whether we mean density or flow direction: \(0\) means density, while a spatial index means flow along the corresponding direction.

So \(T^{00}\) is energy density. The entries \(T^{0i}\) are the three flows of energy. The entries \(T^{i0}\) are the three momentum densities. The entries \(T^{ij}\) are the flows of momentum, and later they will also be recognized as the stress or pressure components.

\begin{figure}[ht]
\centering
\includegraphics[width=0.82\textwidth]{lecture_09_figure_02.png}
\caption{Energy-momentum tensor entries and continuity law}
\end{figure}

The board is partially occluded, so the full tensor must be written as a cautious standard completion:
\begin{equation}
T^{\mu\nu}=
\begin{pmatrix}
T^{00} & T^{01} & T^{02} & T^{03}\\
T^{10} & T^{11} & T^{12} & T^{13}\\
T^{20} & T^{21} & T^{22} & T^{23}\\
T^{30} & T^{31} & T^{32} & T^{33}
\end{pmatrix}.
\end{equation}

A nontrivial fact, emphasized by Susskind, is that for relativistic systems this tensor is symmetric:
\begin{equation}
T^{\mu\nu}=T^{\nu\mu}.
\end{equation}
This is why momentum density and energy flux are equal entries viewed from opposite sides of the diagonal.

For energy itself we have a continuity equation
\begin{equation}
D_\mu T^{0\mu}=0.
\end{equation}
The same argument applies to each momentum component, so altogether we obtain
\begin{equation}
D_\mu T^{\nu\mu}=0.
\end{equation}
Since \(T^{\mu\nu}\) is symmetric, we may also write \(D_\mu T^{\mu\nu}=0\).

This is the real lesson of the long middle part of the lecture. Energy density alone is incomplete. The right-hand side of the gravitational field equation must be the full energy-momentum tensor.

\section{Why the Field Equation Must Be a Tensor Equation}

Once the source is tensorial, the old idea of writing down only a \(00\)-equation becomes impossible. Susskind pauses here to make the covariance argument slowly.

Suppose one were to claim that a law of physics said only \(A_3=B_3\) for two vectors. If that law were to hold in every frame, then by rotating the axes the third component could be turned into the second or the first. The law would then force \(A_2=B_2\) and \(A_1=B_1\) as well. A law that singles out one component cannot be a genuine frame-independent law unless it is really part of a full vector equation.

The same logic applies to relativity, now including time. If a law is to be true in every frame, it must be a tensor equation. A statement about just one preferred component cannot survive Lorentz transformations unless all the corresponding components agree.

So the relativistic generalization of the Newtonian source term cannot stop at the energy density \(T_{00}\). The right-hand side must be the full symmetric tensor
\begin{equation}
8\pi G\,T_{\mu\nu}.
\end{equation}
Accordingly, the left-hand side must also be a symmetric rank-two tensor, built out of the metric and its derivatives. Following the lecture, we may temporarily call it \(G_{\mu\nu}\), not yet because we know what it is, but because we know the kind of thing it must be:
\begin{equation}
G_{\mu\nu}=8\pi G\,T_{\mu\nu}.
\end{equation}

At this point the question becomes precise. What rank-two tensor can be made from the metric, with two derivatives, and with the right conservation property?

\section{Geometric Candidates Built from the Metric}

The left-hand side must be geometric, not material. It must come from the metric. And to resemble the Newtonian second-derivative structure, it should involve two derivatives of the metric.

Susskind only wants reminders here, not a marathon calculation, so the formulas below are best understood as cautious standard reconstructions of the objects on the board. The Christoffel symbols contain first derivatives of the metric:
\begin{equation}
\Gamma^\sigma_{\tau\nu}
=
\frac12 g^{\sigma\delta}
\left(
\partial_\tau g_{\delta\nu}
+
\partial_\nu g_{\delta\tau}
-
\partial_\delta g_{\nu\tau}
\right).
\end{equation}
The Riemann tensor then contains derivatives of the Christoffel symbols, together with terms quadratic in them:
\begin{equation}
R^\mu{}_{\nu\tau\sigma}
=
\partial_\tau\Gamma^\mu_{\nu\sigma}
-
\partial_\sigma\Gamma^\mu_{\nu\tau}
+
\Gamma^\mu_{\lambda\tau}\Gamma^\lambda_{\nu\sigma}
-
\Gamma^\mu_{\lambda\sigma}\Gamma^\lambda_{\nu\tau}.
\end{equation}
The important point is structural: the first two terms contain second derivatives of the metric, while the last two contain squares of first derivatives.

But the full Riemann tensor has four indices, while the source tensor has two. So we contract. The contraction produces the Ricci tensor,
\begin{equation}
R_{\mu\nu}=R^\lambda{}_{\mu\lambda\nu},
\end{equation}
and contracting once more gives the curvature scalar,
\begin{equation}
R=g^{\mu\nu}R_{\mu\nu}.
\end{equation}

\begin{figure}[ht]
\centering
\includegraphics[width=0.78\textwidth]{lecture_09_figure_03.png}
\caption{Candidate geometric tensors for Einstein's equation}
\end{figure}

This produces exactly the exploratory board layout that drives the next stage of the lecture:
\begin{align}
G^{\mu\nu} &= 8\pi G\,T^{\mu\nu},\\
R^{\mu\nu} &\stackrel{?}{=} 8\pi G\,T^{\mu\nu},\\
g^{\mu\nu}R &\stackrel{?}{=} 8\pi G\,T^{\mu\nu}.
\end{align}
The lower two lines are not yet wrong. They are the natural candidates. Each has the right index structure, and each is built from the metric with the right sort of derivatives. So which one do we choose?

We choose by using one more fact.

\section{The Divergence Test}

Susskind now introduces the decisive extra principle with a phrase he likes: we know one more thing. What we know is the continuity equation for energy and momentum,
\begin{equation}
D_\mu T^{\mu\nu}=0.
\end{equation}
If the right-hand side has vanishing covariant divergence, then the left-hand side must as well. Otherwise the equation cannot possibly be right unless we abandon local conservation.

\begin{figure}[ht]
\centering
\includegraphics[width=0.84\textwidth]{lecture_09_figure_04.png}
\caption{Product rule for \(D_\mu(g^{\mu\nu}R)\)}
\end{figure}

\subsection{Question \& Answer}

\paragraph{Question.} Why do the naive candidates \(R^{\mu\nu}\) and \(g^{\mu\nu}R\) fail, and why does the Einstein combination work?

\paragraph{Answer.} Because neither candidate is covariantly conserved by itself.

Let us test the scalar-curvature candidate first. The intended board step is the product rule
\begin{equation}
D_\mu\!\left(g^{\mu\nu}R\right)
=
\left(D_\mu g^{\mu\nu}\right)R
+
g^{\mu\nu}\left(D_\mu R\right).
\end{equation}
Now
\begin{equation}
D_\mu g^{\mu\nu}=0,
\end{equation}
not because of continuity, but because of metric compatibility. And since \(R\) is a scalar,
\begin{equation}
D_\mu R=\partial_\mu R.
\end{equation}
Therefore
\begin{equation}
D_\mu\!\left(g^{\mu\nu}R\right)=g^{\mu\nu}\partial_\mu R.
\end{equation}
In general this is not zero. No good.

What about the Ricci tensor? The lecture quotes the corresponding result in board-level form:
\begin{equation}
D_\mu R^{\mu\nu}
=
\frac12 g^{\mu\nu}D_\mu R
=
\frac12 D^\nu R.
\end{equation}
Again, this is not zero in general. So that candidate also fails.

But now the cure is visible. The divergence of \(R^{\mu\nu}\) gives one half of the derivative term, while the divergence of \(g^{\mu\nu}R\) gives the full derivative term. The combination
\begin{equation}
R^{\mu\nu}-\frac12 g^{\mu\nu}R
\end{equation}
therefore has vanishing covariant divergence. Lowering indices, we define the Einstein tensor by
\begin{equation}
G_{\mu\nu}=R_{\mu\nu}-\frac12 g_{\mu\nu}R,
\end{equation}
and this satisfies
\begin{equation}
D_\mu G^{\mu\nu}=0.
\end{equation}

This is the mathematical climax of the lecture. Among symmetric rank-two tensors built from the metric and two derivatives, there is nothing else with the required conservation property except an overall multiple of this one. Susskind also remarks that there is a more elegant action-principle derivation, but this first route is the one he wants us to see: guess, test, reject, combine.

\section{Einstein's Equation, Vacuum, and the Trace}

Having found the left-hand side, we can finally write the field equation:
\begin{equation}
G_{\mu\nu}=8\pi G\,T_{\mu\nu},
\end{equation}
that is,
\begin{equation}
R_{\mu\nu}-\frac12 g_{\mu\nu}R=8\pi G\,T_{\mu\nu}.
\end{equation}
Susskind remarks, cautiously, that Einstein seems first to have tried the naive Ricci-tensor equation and then discovered that it did not work. The extra \(-\frac12 g_{\mu\nu}R\) term is the crucial correction.

The coefficient is fixed by the Newtonian limit. In a weak field, with slowly moving sources, \(T_{00}\) is essentially \(\rho\), the energy density dominates all the other components, and the geometric side reduces to the same second-derivative texture suggested earlier by \(\nabla^2 g_{00}\sim 8\pi G\rho\). The lecture reports the outcome rather than deriving it in detail: the correct coefficient is exactly the one displayed above.

This also tells us something important about the source. Gravity is not sourced only by energy density. In principle it is sourced by the whole energy-momentum tensor: energy flux, momentum density, and momentum flux as well. In ordinary nonrelativistic situations those other components are typically much smaller, which is why \(T_{00}\) dominates and Newtonian gravity works so well.

\paragraph{Vacuum first.} Susskind next turns to the case with no sources. If
\begin{equation}
T_{\mu\nu}=0,
\end{equation}
then Einstein's equation becomes
\begin{equation}
R_{\mu\nu}-\frac12 g_{\mu\nu}R=0.
\end{equation}
Equivalently,
\begin{equation}
R_{\mu\nu}=\frac12 g_{\mu\nu}R.
\end{equation}
Now contract the free indices. Using \(g^\mu{}_\nu=\delta^\mu{}_\nu\) and \(\delta^\mu{}_\mu=4\), we get
\begin{equation}
R=\frac12\,\delta^\mu{}_\mu\,R=\frac12(4)R=2R.
\end{equation}
Hence
\begin{equation}
R=0,
\end{equation}
and therefore
\begin{equation}
R_{\mu\nu}=0.
\end{equation}
So in vacuum the field equations simplify to the Ricci-flat equation \(R_{\mu\nu}=0\).

\subsection{Question \& Answer}

\paragraph{Question.} If the vacuum equation is \(R_{\mu\nu}=0\), why is vacuum spacetime not automatically flat?

\paragraph{Answer.} Because the Ricci tensor is only a contraction of the full Riemann tensor. It carries less information than the full curvature. The full curvature can remain nonzero even when the Ricci tensor vanishes.

This is one of the lecture's most important payoffs. Schwarzschild spacetime outside its singularity satisfies the vacuum equation, so it is Ricci-flat, but it certainly is not flat. There are tidal forces and genuine curvature there. Gravitational waves are another example. They satisfy \(R_{\mu\nu}=0\) with no source present, just as electromagnetic waves can propagate in regions without charges, but they are most certainly not flat spacetime.

So Ricci-flat is not the same as flat. The contraction from Riemann to Ricci throws away real information.

\paragraph{An audience question pushes one step farther.} After the vacuum discussion an audience member asks what happens if we perform the same contraction without setting \(T_{\mu\nu}\) to zero. This produces the trace relation and the Ricci-form rewriting of Einstein's equation.

\begin{figure}[ht]
\centering
\includegraphics[width=0.82\textwidth]{lecture_09_figure_05.png}
\caption{Trace relation and Ricci-form Einstein equation}
\end{figure}

Define the trace of the energy-momentum tensor by
\begin{equation}
T=g^{\mu\nu}T_{\mu\nu}.
\end{equation}
Now contract Einstein's equation:
\begin{align}
g^{\mu\nu}R_{\mu\nu}
-
\frac12 g^{\mu\nu}g_{\mu\nu}R
&=
8\pi G\,g^{\mu\nu}T_{\mu\nu},\\
R-\frac12(4)R&=8\pi G\,T.
\end{align}
Therefore
\begin{equation}
R=-8\pi G\,T.
\end{equation}
Substituting this back into Einstein's equation yields
\begin{equation}
R_{\mu\nu}=8\pi G\left(T_{\mu\nu}-\frac12 g_{\mu\nu}T\right).
\end{equation}
So one may write the field equation either in Einstein-tensor form or in Ricci form, provided the right-hand side is correspondingly adjusted.

Susskind briefly comments on the trace \(T\). It is zero for electromagnetic radiation and, more generally, for massless particles such as photons or gravitons. For matter with mass it need not vanish.

\paragraph{Two honest closing remarks.} First, Susskind does not claim to have a simple transparent physical interpretation of the Ricci tensor or Ricci scalar comparable to the intuitive meaning of the full Riemann tensor. The Riemann tensor has a clearer geometric role; the contracted objects are simpler, but their individual meaning is less transparent.

Second, the cosmological constant can be treated from either side of the equation. One may regard it as an effective stress-energy contribution
\begin{equation}
T_{\mu\nu}^{(\Lambda)}=\lambda g_{\mu\nu},
\end{equation}
or move the corresponding term to the geometric side and think of it as part of the geometry. From the viewpoint of the equation itself, either bookkeeping convention is possible.

\section{Summary}

The lecture builds Einstein's equation in the same order that one would discover it if one were forced to think structurally rather than computationally. Newton provides the model: field tells matter how to move, matter tells the field what to be. The weak-field Schwarzschild comparison then suggests, only heuristically, that matter should source second derivatives of the metric. But once relativity is taken seriously, the source cannot be a scalar density alone. It must be the full symmetric energy-momentum tensor \(T_{\mu\nu}\), obeying the continuity equation \(D_\mu T^{\mu\nu}=0\).

That continuity equation becomes the selection principle on the geometric side. Neither \(R_{\mu\nu}\) nor \(g_{\mu\nu}R\) is covariantly conserved by itself. The special combination
\begin{equation}
G_{\mu\nu}=R_{\mu\nu}-\frac12 g_{\mu\nu}R
\end{equation}
is. Matching to the Newtonian limit fixes the coefficient, and the field equation becomes
\begin{equation}
G_{\mu\nu}=8\pi G\,T_{\mu\nu}.
\end{equation}
In vacuum this reduces to \(R_{\mu\nu}=0\), but that is only Ricci-flatness, not flatness. Schwarzschild exterior geometry and gravitational waves remain as the lecture's standing reminders that empty spacetime can still be richly curved.